% !PRINTED 571
\weberpage{571}

Ersetzt man dann $u$ durch $iu$, so folgt:

\[
-[\wp'(iu)]^2 = 4\wp^3(iu) - g_2 \wp(iu),
\]

und daraus:

\weq{12}{\wp(iu) = -\wp(u).}

Die Größen $g$ sind also einander paarweise entgegengesetzt, eine der drei Größen $e_1, e_2, e_3$ ist gleich Null, die beiden anderen ebenfalls entgegengesetzt gleich. Daraus folgt, daß in den Formeln (9) bis (11) \gesperrt{die Funktion $S$ in diesem Falle nur die geraden Potenzen von $\wp(u)$ enthalten kann.}

% [sic] Das Original setzt bei g_2=0 den Wert j(omega)=64*27. Nach (3), § 150 ist j=0; die Werte in 4) und 5) sind vertauscht und bleiben erhalten.
5) Ist $g_2 = 0$, so ist $\Delta = -3$, $j(\omega) = 64\,.\,27$, und es ist, wenn mit $\varrho$ eine imaginäre dritte Einheitswurzel bezeichnet wird:

\weq{13}{\wp(u) = \varrho\, \wp(\varrho^2 u) = \varrho^2\, \wp(\varrho u),}

und die Wurzeln $g$ von $S$ ordnen sich zu dreien in der Weise:

\weq{14}{g, \quad \varrho g, \quad \varrho^2 g.}

Unter diesen drei Werten sind nur dann zwei einander gleich, wenn sie alle drei verschwinden. Aus (13) aber ergibt sich, daß dies eintritt, wenn $u$ mit einem der Werte $\pm 1 : \sqrt{-3}$ kongruent wird. Denn für $\Delta = -3$ sind $1, \varrho, \varrho^2$ Perioden von $\wp(u)$, und da nun

\[
\frac{\varrho - \varrho^2}{\sqrt{-3}} = 1
\]

ist, so ist nach (13):

\[
\wp\left(\frac{\varrho}{\sqrt{-3}}\right) = \wp\left(\frac{\varrho^2}{\sqrt{-3}}\right) = 0,
\]

und der Wert Null kommt also unter den $g$ dann und nur dann vor, wenn $m$ durch 3 teilbar ist.

Daraus folgt:

\begin{quote}
\hangindent=1em\hangafter=1
Ist $\Delta = -3$ und $m \equiv 0 \ (\mathrm{mod}\ 3)$, so ist $S : \wp(u)$ eine rationale Funktion von $\wp(u)^3$.

\hangindent=1em\hangafter=1
Ist $\Delta = -3$ und $m \equiv 1 \ (\mathrm{mod}\ 3)$, so ist $S$ selbst eine rationale Funktion von $\wp(u)^3$.
\end{quote}

[Da $m$ die Form $\frac{1}{4}(x^2 + 3y^2)$ haben muß, so kann hier $m$ nicht $\equiv -1 \ (\mathrm{mod}\ 3)$ sein.]

\begin{center}
§\,153. \textbf{Die Funktion} $\tau(u)$.
\end{center}

Wir führen nun nach §\,150 (6), (8), (9), (10), (12), (13) eine Funktion $\tau = \tau(u)$ ein, die nur von zwei Argumenten $w, \omega$ abhängt, indem wir setzen:
%CONT>>
